\documentclass[11pt]{book}

\usepackage{hyperref}
\usepackage{listings}
\usepackage{minted}
\usepackage{tcolorbox}
\usepackage{fontspec}
\usepackage{geometry}
\usepackage{fancyhdr}
\usepackage{xcolor}

% Custom colors for different prompt components
\definecolor{promptcolor}{RGB}{50,120,190}
\definecolor{contextcolor}{RGB}{70,150,70}
\definecolor{constraintcolor}{RGB}{180,90,90}

% Custom environments for prompt engineering
\newenvironment{prompt}[1]
  {\begin{tcolorbox}[colback=promptcolor!5,colframe=promptcolor!50,title=Prompt: #1]}
  {\end{tcolorbox}}

\newenvironment{context}
  {\begin{tcolorbox}[colback=contextcolor!5,colframe=contextcolor!50,title=Context]}
  {\end{tcolorbox}}

\newenvironment{constraints}
  {\begin{tcolorbox}[colback=constraintcolor!5,colframe=constraintcolor!50,title=Constraints]}
  {\end{tcolorbox}}

\title{PACTS Promptbase Engineering}
\author{Project Documentation}
\date{\today}

\begin{document}

\maketitle
\tableofcontents

\chapter{Prompt Engineering Principles}

\section{Core Principles}
The PACTS promptbase follows these core principles:
\begin{itemize}
  \item Context-aware prompting
  \item Safety-first design
  \item Verifiable outputs
  \item Transparent reasoning
\end{itemize}

\section{Prompt Components}
Each prompt in PACTS consists of:
\begin{itemize}
  \item Context block
  \item Instruction block
  \item Constraints block
  \item Verification block
\end{itemize}

\chapter{Base Prompts}

\section{Permission Management}
\begin{prompt}{Permission Check}
\begin{minted}{dhall}
let PermissionCheck =
      { Type =
          { action : Text
          , context : List Text
          , constraints : List Text
          , verification : List Text
          }
      , default =
          { action = ""
          , context = [] : List Text
          , constraints = [] : List Text
          , verification = [] : List Text
          }
      }
\end{minted}

\begin{context}
You are a permission verification system. Your task is to evaluate whether the requested action complies with system policies and safety constraints.

Current system state:
\begin{minted}{json}
{
  "mode": "strict",
  "safetyLevel": "high",
  "activeConstraints": ["AI_SAFETY", "DATA_PRIVACY"]
}
\end{minted}
\end{context}

\begin{constraints}
1. All actions must be explicitly permitted
2. Safety constraints cannot be overridden
3. Context must be fully validated
4. Decisions must be explained
\end{constraints}
\end{prompt}

\section{Compliance Verification}
\begin{prompt}{Compliance Check}
\begin{minted}{dhall}
let ComplianceCheck =
      { Type =
          { rules : List Text
          , evidence : List Text
          , requirements : List Text
          , validation : List Text
          }
      , default =
          { rules = [] : List Text
          , evidence = [] : List Text
          , requirements = [] : List Text
          , validation = [] : List Text
          }
      }
\end{minted}

\begin{context}
You are a compliance verification system. Your task is to evaluate whether the provided evidence meets the specified compliance requirements.

Compliance framework:
\begin{minted}{json}
{
  "framework": "AI_SAFETY",
  "version": "1.0",
  "requirements": [
    "explainability",
    "reproducibility",
    "safety_bounds"
  ]
}
\end{minted}
\end{context}

\begin{constraints}
1. All requirements must be verified
2. Evidence must be documented
3. Validation steps must be traceable
4. Non-compliance must be reported
\end{constraints}
\end{prompt}

\chapter{Prompt Composition}

\section{Combining Prompts}
\begin{minted}{dhall}
let CombinedCheck =
      λ(pc : Type) →
      λ(cc : Type) →
        { permission = pc
        , compliance = cc
        , combined = { permission : pc, compliance : cc }
        }
\end{minted}

\section{Prompt Inheritance}
\begin{minted}{dhall}
let BasePrompt =
      { Type =
          { context : List Text
          , constraints : List Text
          }
      }

let SpecializedPrompt =
      λ(base : Type) →
        { inherit = base
        , additional = List Text
        }
\end{minted}

\chapter{Verification Patterns}

\section{Safety Verification}
\begin{prompt}{Safety Check}
\begin{minted}{dhall}
let SafetyCheck =
      { Type =
          { bounds : List Text
          , assertions : List Text
          , invariants : List Text
          }
      , default =
          { bounds = [] : List Text
          , assertions = [] : List Text
          , invariants = [] : List Text
          }
      }
\end{minted}

\begin{context}
You are a safety verification system. Your task is to ensure that all operations remain within defined safety bounds.
\end{context}

\begin{constraints}
1. All safety bounds must be checked
2. Invariants must be maintained
3. Violations must trigger immediate stops
\end{constraints}
\end{prompt}

\chapter{Implementation}

\section{Generating Prompts}
\begin{minted}{typescript}
import { dhall } from '@dhall/core';
import { PromptTemplate } from '@pacts/core';

export async function generatePrompt(
  template: string,
  context: Record<string, unknown>
): Promise<string> {
  const config = await dhall.load('./prompts/config.dhall');
  const promptTemplate = new PromptTemplate(template, config);
  return promptTemplate.render(context);
}
\end{minted}

\section{Validating Prompts}
\begin{minted}{typescript}
import { z } from 'zod';
import { PromptValidator } from '@pacts/core';

const PromptSchema = z.object({
  context: z.array(z.string()),
  constraints: z.array(z.string()),
  verification: z.array(z.string())
});

export async function validatePrompt(
  prompt: unknown
): Promise<boolean> {
  return PromptValidator.validate(prompt, PromptSchema);
}
\end{minted}

\chapter{Testing}

\section{Prompt Testing Framework}
\begin{minted}{typescript}
import { test } from '@jest/globals';
import { PromptTester } from '@pacts/core';

test('Permission Check Prompt', async () => {
  const tester = new PromptTester();
  const result = await tester.test('permission_check', {
    action: 'read_file',
    context: { user: 'admin' }
  });
  expect(result.isValid).toBe(true);
});
\end{minted}

\end{document} 